% GENERATED FILE. Edit canonical lessons, then run npm run generate:books.
% canonical-source-hash: fnv1a64:1f93622ac4667c86
% canonical-lessons: IT-C467-lampo, IT-C467-tuonare, IT-C467-lampeggiare, IT-C467-grandine, IT-C467-grandinare

\chapter{Il temporale}
\label{ch:it-il-temporale}
\hlchaptermodality{467}

\begin{chapteropening}
\textbf{By the end of this chapter:} \emph{I can describe a thunderstorm: lightning, thunder, flashing, hail and hailing.}

Complete the last lesson of chapter 467: I can describe a thunderstorm: lightning, thunder, flashing, hail and hailing.
\end{chapteropening}

\section[il lampo]{il lampo — the flash of lightning, the flash}

\label{lesson:IT-C467-lampo}



\begin{quote}
Before the new one: say the Italian for the drought, then the Italian for the breeze.
\end{quote}

\subsection*{You'll want to know: il lampo}
\textbf{il lampo} — \textquotedblleft{}the flash of lightning, the flash\textquotedblright{}.

\begin{quote}
Ho visto un lampo sopra le colline.
\end{quote}

\begin{quote}
\emph{I saw a flash of lightning over the hills.}
\end{quote}

\emph{In un lampo} — in a flash.

\begin{cousinweb}
\emph{lampo} goes with \emph{lampa} 'lamp', which EtymDB traces through Old French \emph{lampe} to Late Latin \emph{lampada}, from Greek \emph{lampás} 'torch'.
\end{cousinweb}

\begin{grammarlens}[title={What you've built}]
One more word for telling, explaining and describing at B1.
\end{grammarlens}

\subsection*{Your turn}
\begin{itemize}
\raggedright
  \item \emph{Say it:} \emph{il lampo}
  \item \emph{Say it:} \emph{il lampo}, once more
  \item \emph{Say it:} a sentence of your own with \emph{il lampo}
  \item \emph{Recall:} say the Italian for the drought, then the Italian for the breeze, then say \emph{il lampo} again
\end{itemize}

\begin{tcolorbox}[breakable,colback=teal!4,colframe=teal!35!black,title={Before you move on}]
What does il lampo mean? (\textquotedblleft{}The flash of lightning, the flash\textquotedblright{}.) Say it once more.
\end{tcolorbox}
\section[tuonare]{tuonare — to thunder}

\label{lesson:IT-C467-tuonare}



\begin{quote}
Before the new one: say the Italian for the breeze, then the Italian for the flash of lightning.
\end{quote}

\subsection*{You'll want to know: tuonare}
\textbf{tuonare} — \textquotedblleft{}to thunder\textquotedblright{}.

\begin{quote}
Tuona da un'ora, ma non piove.
\end{quote}

\begin{quote}
\emph{It has been thundering for an hour, but it is not raining.}
\end{quote}

\emph{Il tuono} — thunder.

\begin{cousinweb}
\emph{tuonare} is inherited from Latin \emph{tonāre} 'to thunder' (EtymDB).
\end{cousinweb}

\begin{grammarlens}[title={What you've built}]
One more word for telling, explaining and describing at B1.
\end{grammarlens}

\subsection*{Your turn}
\begin{itemize}
\raggedright
  \item \emph{Say it:} \emph{tuonare}
  \item \emph{Say it:} \emph{tuonare}, once more
  \item \emph{Say it:} a sentence of your own with \emph{tuonare}
  \item \emph{Recall:} say the Italian for the breeze, then the Italian for the flash of lightning, then say \emph{tuonare} again
\end{itemize}

\begin{tcolorbox}[breakable,colback=teal!4,colframe=teal!35!black,title={Before you move on}]
What does tuonare mean? (\textquotedblleft{}To thunder\textquotedblright{}.) Say it once more.
\end{tcolorbox}
\section[lampeggiare]{lampeggiare — to flash (with lightning); to blink}

\label{lesson:IT-C467-lampeggiare}



\begin{quote}
Before the new one: say the Italian for the flash of lightning, then the Italian for to thunder.
\end{quote}

\subsection*{You'll want to know: lampeggiare}
\textbf{lampeggiare} — \textquotedblleft{}to flash (with lightning); to blink\textquotedblright{}.

\begin{quote}
Lampeggia lontano, sul mare.
\end{quote}

\begin{quote}
\emph{Lightning is flashing far off, over the sea.}
\end{quote}

In a car: \emph{la freccia lampeggia} — the indicator is flashing.

\begin{cousinweb}
\emph{lampeggiare} is built on \emph{lampo} 'flash', which goes with \emph{lampa} 'lamp'; EtymDB traces \emph{lampa} through Old French \emph{lampe} to Late Latin \emph{lampada}, from Greek \emph{lampás} 'torch'.
\end{cousinweb}

\begin{grammarlens}[title={What you've built}]
One more word for telling, explaining and describing at B1.
\end{grammarlens}

\subsection*{Your turn}
\begin{itemize}
\raggedright
  \item \emph{Say it:} \emph{lampeggiare}
  \item \emph{Say it:} \emph{lampeggiare}, once more
  \item \emph{Say it:} a sentence of your own with \emph{lampeggiare}
  \item \emph{Recall:} say the Italian for the flash of lightning, then the Italian for to thunder, then say \emph{lampeggiare} again
\end{itemize}

\begin{tcolorbox}[breakable,colback=teal!4,colframe=teal!35!black,title={Before you move on}]
What does lampeggiare mean? (\textquotedblleft{}To flash (with lightning); to blink\textquotedblright{}.) Say it once more.
\end{tcolorbox}
\section[la grandine]{la grandine — the hail}

\label{lesson:IT-C467-grandine}



\begin{quote}
Before the new one: say the Italian for to thunder, then the Italian for to flash.
\end{quote}

\subsection*{You'll want to know: la grandine}
\textbf{la grandine} — \textquotedblleft{}the hail\textquotedblright{}.

\begin{quote}
La grandine ha rovinato le macchine.
\end{quote}

\begin{quote}
\emph{The hail damaged the cars.}
\end{quote}

\emph{Un chicco di grandine} — a hailstone.

\begin{cousinweb}
\emph{grandine} is inherited from Latin \emph{grandō, grandinem} 'hail' (EtymDB).
\end{cousinweb}

\begin{grammarlens}[title={What you've built}]
One more word for telling, explaining and describing at B1.
\end{grammarlens}

\subsection*{Your turn}
\begin{itemize}
\raggedright
  \item \emph{Say it:} \emph{la grandine}
  \item \emph{Say it:} \emph{la grandine}, once more
  \item \emph{Say it:} a sentence of your own with \emph{la grandine}
  \item \emph{Recall:} say the Italian for to thunder, then the Italian for to flash, then say \emph{la grandine} again
\end{itemize}

\begin{tcolorbox}[breakable,colback=teal!4,colframe=teal!35!black,title={Before you move on}]
What does la grandine mean? (\textquotedblleft{}The hail\textquotedblright{}.) Say it once more.
\end{tcolorbox}
\section[grandinare]{grandinare — to hail}

\label{lesson:IT-C467-grandinare}



\begin{quote}
Before the new one: say the Italian for to flash, then the Italian for the hail.
\end{quote}

\subsection*{You'll want to know: grandinare}
\textbf{grandinare} — \textquotedblleft{}to hail\textquotedblright{}.

\begin{quote}
Ieri ha grandinato per dieci minuti.
\end{quote}

\begin{quote}
\emph{Yesterday it hailed for ten minutes.}
\end{quote}

It takes \emph{avere} or \emph{essere}: \emph{ha grandinato}, \emph{è grandinato}.

\begin{cousinweb}
\emph{grandinare} is inherited from Latin \emph{grandināre} 'to hail' (EtymDB), from \emph{grandō} 'hail'.
\end{cousinweb}

\begin{grammarlens}[title={What you've built}]
One more word for telling, explaining and describing at B1.
\end{grammarlens}

\subsection*{Your turn}
\begin{itemize}
\raggedright
  \item \emph{Say it:} \emph{grandinare}
  \item \emph{Say it:} \emph{grandinare}, once more
  \item \emph{Say it:} a sentence of your own with \emph{grandinare}
  \item \emph{Recall:} say the Italian for to flash, then the Italian for the hail, then say \emph{grandinare} again
\end{itemize}

\begin{tcolorbox}[breakable,colback=teal!4,colframe=teal!35!black,title={Before you move on}]
What does grandinare mean? (\textquotedblleft{}To hail\textquotedblright{}.) Say it once more.
\end{tcolorbox}
